\subsection{Computational complexity of exact compression}
\label{subsec:aor-safe-compression-complexity}

We use an explicit finite encoding. A source instance lists a finite strategy
library $L=\{s_0,s_1,\ldots,s_n\}$, a finite nonempty belief set
$\mathcal B$, exact binary-rational profiles $j_s(b)$, a finite module
universe $\mathcal M$, and the incidence rows
$\mods(s)\subseteq\mathcal M$. Strategy $s_0$ is the mandatory inactive
strategy, has profile zero at every belief, carries no modules, and has zero
burden. Every active strategy has a positive binary-integer weight $w_i$, and
the threshold $B$ is a nonnegative binary integer. Closure is identity
closure,

\[
  C(A)=\bigcup_{s\in A}\mods(s).
\]

The source frontier is $F_L(b)=\max_{s\in L}j_s(b)$, and the burden of a
sublibrary is $W(A)=\sum_{s_i\in A,\,i\geq 1}w_i$. The decision problem
\textsc{Identity-Closure Safe Compression} (\textsc{ISC-ID}) asks whether
there is an $A\subseteq L$ such that $s_0\in A$, $W(A)\leq B$,
$F_A=F_L$ at every listed belief, and $C(A)=C(L)$.

The hardness source below is a deliberately normalized version of weighted
set cover. In \textsc{Full-Union Weighted Set Cover} (\textsc{FU-WSC}), the
input is a nonempty finite universe $U$, an indexed finite family
$(A_i)_{i=1}^n$ whose union is $U$, positive binary-integer costs $c_i$, and a
nonnegative binary-integer budget $B$. The question is whether some index set
$I\subseteq\{1,\ldots,n\}$ satisfies
$\bigcup_{i\in I}A_i=U$ and $\sum_{i\in I}c_i\leq B$.
Encodings that fail these stated syntactic restrictions are rejected, so
\textsc{FU-WSC} is a language rather than an unverified promise problem.

\begin{lemma}[Normalized weighted set cover]
\label{lem:aor-full-union-weighted-set-cover}
\textnormal{\textsc{FU-WSC}} is NP-complete.
\end{lemma}

\begin{proof}
Input validity is checkable in polynomial time. A selected-index bit vector is
a polynomial certificate: unioning its listed incidence rows and adding its
binary-integer costs verifies feasibility in polynomial time. Thus
\textsc{FU-WSC} is in NP.

For hardness, reduce the standard \textsc{Set Cover} decision problem, one of
Karp's NP-complete covering problems \citep{Karp1972}. Write an input as
$(U,(A_i)_{i=1}^n,k)$, with nonnegative cardinality bound $k$. If
$U=\varnothing$, output the fixed yes-instance consisting of universe
$\{u\}$, the one set $\{u\}$, cost one, and budget one. If
$U\neq\varnothing$ but $\bigcup_i A_i\neq U$, output the fixed no-instance
with universe $\{u\}$, the one set $\{u\}$, cost two, and budget one.
Otherwise retain the universe and family, assign every set unit cost, and set
the budget to $k$. Every output has a nonempty universe, a full-union family,
positive integer costs, and a nonnegative integer budget. The first branch
preserves the fact that the empty universe is coverable; the second preserves
the fact that an uncovered element makes the source instance a no-instance;
and in the third branch cost is exactly selected-set cardinality. The case
split, union check, and copied incidence representation have polynomial size
and construction time. Hence the mapping preserves yes- and no-instances and
proves NP-hardness.
\end{proof}

\begin{theorem}[Identity-closure safe-compression complexity]
\label{thm:aor-safe-compression-complexity}
\textnormal{\textsc{ISC-ID}} is NP-complete. More specifically, it remains
NP-complete on
the subclass with exactly one belief, zero profiles for every active strategy,
an inactive strategy with zero profile, zero burden, and no modules, and a
source library whose active module rows have full union. The result holds even
when every active weight is one. Consequently, minimum-burden identity-closure
safe compression is NP-hard. Any broader safe-compression model class
containing this identity-closure subclass is NP-hard; no membership claim for
an arbitrary nonidentity closure operator follows without a polynomial
closure-equality verifier in its declared encoding.
\end{theorem}

\begin{proof}
\emph{Membership in NP.}
Encode a proposed retained library by one bit per source strategy. Inactive
retention and source inclusion are direct bit checks. Binary-integer addition
checks $W(A)\leq B$. For every explicitly listed belief, exact rational
comparison finds the source and candidate maxima and checks their equality.
Under identity closure, bitwise union of the listed module-incidence rows
checks $C(A)=C(L)$. Integer and rational operations have bit complexity
polynomial in the explicit input length. Thus \textsc{ISC-ID} is in NP.

\emph{Reduction algorithm.}
Given a valid \textsc{FU-WSC} instance
$(U,(A_i)_{i=1}^n,(c_i)_{i=1}^n,B)$, create one module $m_u$ for each
$u\in U$, one active strategy $s_i$ for each set $A_i$, and one mandatory
inactive strategy $s_0$. Use one belief $b_0$. Set

\[
  j_{s_i}(b_0)=j_{s_0}(b_0)=0,
  \qquad
  \mods(s_i)=\{m_u:u\in A_i\},
  \qquad
  \mods(s_0)=\varnothing.
\]

Give $s_i$ weight $c_i$, give $s_0$ weight zero, take the complete constructed
library as the source, use identity closure, and copy the budget without
alteration.

\emph{Polynomial construction size.}
The construction adds one module per universe element, one active strategy per
set, one inactive strategy, and one belief. Its module-incidence table is the
set-cover incidence table; the remaining profiles are $n+1$ zero entries, and
costs and budget are copied. A dense encoding has size $O(n|U|)$ plus the
copied identifier and integer bit lengths, while a sparse encoding is linear
in the copied incidence list plus this overhead. Either encoding is polynomial
in the corresponding explicit source encoding.

\emph{Forward feasibility.}
Suppose $I$ is a set cover with $\sum_{i\in I}c_i\leq B$. Retain
$A=\{s_0\}\cup\{s_i:i\in I\}$. Both source and retained frontiers equal zero
at $b_0$, with $s_0$ itself attaining that value. Identity closure gives

\[
 C(A)=\bigcup_{i\in I}\mods(s_i)
     =\{m_u:u\in U\}
     =\bigcup_{i=1}^n\mods(s_i)=C(L),
\]

where the last equality uses the full-union condition. The inactive strategy
is retained, and $W(A)=\sum_{i\in I}c_i\leq B$, so $A$ is a yes-certificate
for the constructed \textsc{ISC-ID} instance.

\emph{Reverse feasibility.}
Suppose $A$ is a safe constructed sublibrary within threshold, and let
$I=\{i:s_i\in A\}$. Exact closure preservation and the empty inactive module
row imply

\[
 \bigcup_{i\in I}\{m_u:u\in A_i\}
   =C(A)=C(L)=\{m_u:u\in U\}.
\]

The module labels are in one-to-one correspondence with $U$, so
$\bigcup_{i\in I}A_i=U$. Thus $I$ is a set cover. Frontier preservation adds
no hidden selection requirement because the mandatory inactive strategy
already attains the common zero frontier.

\emph{Threshold preservation.}
The inactive strategy has zero burden, active weights are copied exactly, and
the threshold is copied exactly. Therefore
$W(A)=\sum_{i\in I}c_i$ for every selected index set $I$, not merely for
optimal selections. Hence $W(A)\leq B$ if and only if the corresponding set
cover has cost at most $B$. The reduction is a polynomial many-one reduction.
Together with membership in NP and
\cref{lem:aor-full-union-weighted-set-cover}, this proves NP-completeness of
the stated restricted subclass and therefore of \textsc{ISC-ID}. The
normalization proof uses unit costs on its nontrivial branch, so the restricted
subclass remains NP-complete with unit active weights. Solving the associated
minimum-burden problem would decide whether its optimum is at most $B$, which
proves NP-hardness of the optimization problem.
\end{proof}


The encoding, normalization, and both directions above establish the human proof. Exact reduction fixtures are supplementary finite checks.
